\documentclass[12pt]{article}
\usepackage[utf8]{inputenc}
\usepackage{amsmath}
\usepackage{amsfonts}
\usepackage{amssymb}
\usepackage{graphicx}
\usepackage{float}
\usepackage{booktabs}
% --- Add these to mimic MS Word ---

% 1. Margins: Sets standard 1-inch margins on all sides
\usepackage[margin=1in]{geometry} 

% 2. Paragraphs: Removes the first-line indent and adds space between paragraphs
\usepackage{parskip} 

% 3. Line Spacing: Makes the lines slightly less cramped
\usepackage{setspace}
\setstretch{1.15} % 1.15 is MS Word's default. Use \onehalfspacing for 1.5 spacing.

\title{Assignment 5}
\author{Shapnil Das}
\begin{document}

\maketitle

\section{Question:}

Simulate $500$ independent Random Walks. Regress the 1st random walk on the other $499$ (one at a time).
Create a histogram of the $500$ resulting $R^2$ values. What percentage of these regressions yielded a t-statistic greater than $1.96$?
If these variables are totally unrelated, why isn't that percentage $5\%$?

\section{Solution:}

First, we will use loops to generate 500 Independent Random Walks for 500 different shocks.
A Random Walk process can be defined as the sum of all current and past shocks.
These shocks are independent and has a non-decaying effect on the series.

\[y_t = y_{t-1} + \varepsilon_t \]

\[ y_t = \sum_{i=1}^{t} \varepsilon_i \quad \text{for} \; y_0 = 0.\]
where $\varepsilon_t$ is a white noise error term.

we have generated 500 independent shock $\varepsilon_i \sim \mathcal{N}(0,1)$ for \; $i = 1,2,...,500$.
And for each shock we have generated 500 independent random walk process $y_t$ for $t = 1,2,...,500$ as the sum of all current and past shocks
for 100 initial sample or observations $T = 100$.

Now we regress the first random walk $y_1$ on the other 499 random walks
one at a time and calculate the $R^2$ and $T-statistic$ values for each regression.
\[ y_1 = \alpha + \beta y_j + u_t \quad \text{for} \; j = 2,3,...,500 \]
Thus we will have 499 sample statistics which can mimic a Monte-Carlo simulation
of the distribution of $R^2$ and $T-statistic$ values when the variables are totally unrelated.

Our assumption about the $R^2$ score is as the random walks are independent, theoretically there should be any relation between between $y_1$ and $y_j$ for $j = 2,3,...,500$.
The expected value of $R^2$ should be close to zero and the expected value of the $T-statistics$ should be significantly close to zero $95\%$ of the time.
considering the rejection region is 1.96 is in absolute value at 5\% significance level, we should expect to not reject the null hypothesis of no relationship between $y_1$ and $y_j$ for only 5\% of the regressions.

\begin{table}[H]
    \centering
    \renewcommand{\arraystretch}{1.2}
    \setlength{\tabcolsep}{18pt}
    \caption{Summary Statistics}
    \vspace{0.15cm}
    \begin{tabular}{lrrr}
        \toprule
        \textbf{Variable} & \textbf{Obs} & \textbf{Mean} & \textbf{Std. Dev.} \\
        \midrule
        Reject            & 499          & $0.731$       & $0.444$            \\
        $R^2$             & 499          & $0.142$       & $0.126$            \\
        \bottomrule
    \end{tabular}
    \label{tab:summary_stats_booktabs}
\end{table}

By Investigating the summary statistics we can see contradictions to our expectations.The T-statistics is rejecting the null hypothesis
of no relationship between $y_1$ and $y_j$ for 73.1\% of the regressions which is much higher than the expected value in 5\% significance level.
And the $R^2$ values shows the independent random walks explains $14.2\%$ of the variation in the first random walk.

This contradictory results are known as spurious regression.

A standard normal distribution of the shock contains a constant mean and constant variance but the shock of a random walk process follows a constant mean
and a time dependent infinite variance where $\varepsilon_i \sim \mathcal{N}(\mu,t\sigma^2_i)$. Random walks have a non-decaying effect of all current and past shocks on the series, which can lead to create artificial relations with other
unrelated random walks.A small shocks can change the trajectory of a random walk but it can never revert to it's original mean and wanders away to new trajectory randomly as time grows.By pure chance, 2 or more random walks can move in the same or opposite direction.So when we regress one on another,
random walks, the regression can falsely identify the correlation among them as causation or relations.Thus they create an artifical goodness of fit $r^2$ and significant t-statistics which leads to spurious regression results.
This is one of the core drawbacks of using random walk non-stationary process in Time Series analysis.

\includegraphics[width=\textwidth,scale=1.2]{./histogram_hw5.png}

By plotting the histogram of the $R^2$ values we can see that most of the values are close to zero but there are some outliers which are far from zero and creates a long right skewed tail in the distribution.
These outliers are the result of spurious regression which creates a significant t-statistics and high values  of $R^2$ even when the variables are totally unrelated.

\includegraphics[width=\textwidth,scale=1.2]{./kdensity_hw5.png}

From the kernel density plot of the T-statitics we can see the contrast comparison of the standard normal distribution of the t-stat
concentrated around zero and the distribution of the t-statistics from the regressions which is more spread out and presenting an uniform distribution,
meaning standard significance tests will falsely tell us the variables are related.
The 2 black dashed line represents the 2 tail rejection area at 5\% significance level at $\pm$1.96. For our empirical distribution of t-statistics,
we can see that a large portion of the distribution falls outside the rejection area which is the result of spurious regression.



\end{document}